\subsection{Sprint 0}

Como \ac{ages} III, reconheço que as \textit{sprints} iniciais são fundamentais, pois nelas definimos as tecnologias que serão utilizadas ao longo de todo o projeto. Em conjunto com os demais integrantes da equipe de \ac{ages} III, analisamos as necessidades específicas da solução --- que envolvia a construção de um mapa 3D interativo --- e concluímos que a melhor abordagem seria utilizar uma biblioteca especializada em renderização tridimensional no \textit{frontend}.

Dessa forma, definimos a \textit{stack} tecnológica da seguinte maneira: React seria utilizado como base da aplicação, enquanto a biblioteca Three.js ficaria responsável pela renderização e manipulação do mapa 3D. Para a estilização dos componentes, optamos pelo uso de \textit{styled-components}, aproveitando a flexibilidade da abordagem CSS-in-JS. No \textit{backend}, como nenhum integrante do grupo possuía experiência avançada, optamos por uma \textit{stack} que já era familiar para a maioria e que utilizasse a mesma linguagem do \textit{frontend} --- por isso, escolhemos o NestJS, \textit{framework} robusto e bem estruturado em TypeScript, similar ao \textit{frontend}, garantindo que a curva de aprendizado fosse a menor possível.

Com as definições tecnológicas finalizadas, iniciamos a criação dos repositórios do projeto. No entanto, enfrentamos um contratempo significativo: o arquiteto responsável pela \ac{ages} precisou se desligar da agência, o que gerou atrasos na criação e estruturação inicial dos repositórios e causou certa apreensão na equipe. Ainda assim, conseguimos avançar e estruturar o ambiente inicial de desenvolvimento de forma colaborativa.

Com o tempo restante após todas essas definições, concentrei meus esforços em aprofundar o conhecimento na biblioteca Three.js, identificada desde o início como um dos pontos mais críticos do projeto. Essa dedicação inicial foi essencial para garantir segurança técnica na construção do mapa 3D nas etapas seguintes.
